\documentclass[11pt]{article}
\usepackage{acl}
\usepackage{times}
\usepackage{latexsym}
\usepackage[T1]{fontenc}
\usepackage[utf8]{inputenc}
\usepackage{microtype}
\usepackage{inconsolata}
\usepackage{amsmath,amssymb,booktabs}

\title{Drift-Aware Memory-Bounded Batching for Irregular GPU Pipelines: Online Footprint Feedback for Memory-Bounded Workloads}
\author{Lee Si Yuan \\
  Independent Researcher \\
  \texttt{siyuanlee.cs@gmail.com}}

\begin{document}
\maketitle
\begin{abstract}
Memory-bounded batching removes reliance on backend occupancy estimates by filling a worker's memory budget with as much work as its resource model permits. That policy assumes the estimated per-item footprint remains representative as an input stream evolves. We propose Drift-Aware Memory-Bounded Batching (DAMB), a runtime controller that tracks per-class footprint estimates and uses a conservative, smoothed estimate to revise safe batch capacities when workload characteristics change. DAMB keeps explicit memory bounds as a hard constraint and leaves kernel code unchanged. We implement a deterministic CPU simulation of a manager/worker pipeline with three workload phases and compare first-phase static calibration against an online smoothed estimate with a fixed safety reserve. In this simplified model, the online policy reaches 1.822 normalized items per time unit versus 1.805 for the static policy (1.009$\times$); this small difference is sensitive to integer batch partitioning and is not evidence of a GPU speedup. Results are simulation-only; no docking code, molecular data, or GPU hardware was used.
\end{abstract}

\section{Introduction}
Irregular GPU pipelines spend time in host-side decisions such as grouping inputs, allocating buffers, and submitting work across workers and streams. Accordi et al.~\citep{accordi2026} study these out-of-kernel decisions in the muDock molecular docking mini-app. Their memory-bounded policy chooses the largest class-specific batch that fits the memory assigned to a worker, avoiding occupancy hints that may differ across backends.

A fixed footprint model is useful when a run has a stable input distribution. Screening workloads can instead change over time. A stale estimate can leave memory unused or lead to an unsafe allocation. This paper studies a bounded feedback loop that observes completed batches and updates the next batch capacity. Our contributions are an online estimate-and-reserve rule, a deterministic CPU simulation, and a proposed drift-triggered controller and GPU evaluation protocol. This is a feasibility study, not a reproduction of the hardware results in~\citep{accordi2026}.

\section{Method: Drift-Aware Memory-Bounded Batching}
\subsection{Online estimate}
For each class $c$, maintain an exponentially smoothed footprint estimate:
\begin{equation}
\hat m_t(c)=\alpha m^{\mathrm{obs}}_t(c)+(1-\alpha)\hat m_{t-1}(c),
\end{equation}
where $m^{\mathrm{obs}}_t(c)$ is measured from the latest completed batch and $\alpha$ is a smoothing parameter. The initial estimate can come from a warm-up sample or a declared conservative bound.

\subsection{Conservative batch capacity}
For worker memory budget $M_{\mathrm{worker}}$ and reserve factor $\rho\in(0,1]$, compute
\begin{equation}
B_t(c)=\max\left(1,\left\lfloor\frac{\rho M_{\mathrm{worker}}}{\hat m_t(c)}\right\rfloor\right).
\end{equation}
The reserve accounts for model error and non-payload allocations and must be calibrated with measured peak allocation. Where a reliable hard upper bound $m_{\max}(c)$ exists, enforce the stronger safety check
\begin{equation}
B_t(c)m_{\max}(c)+M_{\mathrm{overhead}}\le M_{\mathrm{worker}}.
\end{equation}
An estimate alone cannot guarantee safety under an unbounded distribution shift.

\subsection{Drift trigger and controller policy}
A production controller should update after a fixed observation window and change capacity only when
\begin{equation}
\frac{|m_{\mathrm{obs}}(c)-\hat m(c)|}{\max(\hat m(c),\epsilon)}>\delta.
\end{equation}
On a trigger, it should lower the next capacity immediately, then raise capacity gradually after sustained lower observations. A known class bound or conservative high percentile should cap the estimate. Without a safe bound, retain the existing safe capacity until enough observations are collected. Controller state is per class and device, avoiding transfer of one backend's calibration to another.

\section{Experimental Setup}
The deterministic Python simulation uses seed 20261106. Each of three phases contains 2,400 items, with light-item proportions 0.90, 0.25, and 0.90. Light items use one normalized memory unit and one work unit; heavy items use eight memory units and six work units. The worker budget is 96 units, four workers execute batches, and each batch adds two normalized launch-cost units.

The static baseline uses initial class footprint estimates and the full memory budget. The online policy uses the same starting estimates, a 0.85 reserve factor, and 0.5 exponential smoothing after each phase. Class-specific jobs are grouped, then batches are greedily assigned to the least-loaded worker. Simulated time is the maximum worker load for each phase, summed across phases. These choices define a scheduling sensitivity model, not a calibrated model of a GPU or muDock. The experiment is available in \texttt{adaptive\_batching\_experiment.py}; its configuration and results are recorded in \texttt{results.json}.

\section{Results}
\begin{table}[t]
\centering
\scriptsize
\caption{CPU simulation results.}
\label{tab:results}
\begin{tabular}{lrrr}
\toprule
Policy & Time & Throughput & Relative \\
\midrule
Static calibration & 3988 & 1.805 & 1.000$\times$ \\
Online, 15\% reserve & 3952 & 1.822 & 1.009$\times$ \\
\bottomrule
\end{tabular}
\end{table}
The online policy differs by less than 1\% in this setup. Its phase throughput is 4.131, 0.860, and 4.131, compared with 4.082, 0.853, and 4.082 for static calibration. Class footprints are constant and the initial estimates already match them, so the workload does not meaningfully exercise estimation drift. Integer batch partitioning and greedy worker assignment can account for the small timing difference. This result does not establish a general benefit or a GPU performance gain.

\section{Discussion and Limitations}
The experiment has two fixed classes, no transfers, allocation fragmentation, kernel execution, external GPU contention, or measurement noise. The 0.9\% difference can arise from discrete batch scheduling rather than useful adaptation. The method is most plausible when per-class footprints drift and the runtime can observe allocations or obtain safe class bounds. It may be ineffective when the base model is accurate, and a conservative reserve can reduce throughput. An average estimate alone may be unsafe for heavy-tailed items; hard bounds and asymmetric updates are important for deployment.

A real evaluation should integrate DAMB into muDock or a comparable pipeline, replay time-ordered workloads with distribution shifts, and compare fixed memory-bounded batching, online mean estimates, and high-percentile estimates with a hard safety bound. It should report throughput, peak memory, allocation failures, batch wait time, worker idle time, and per-phase results across repeated runs on each supported backend. No such GPU evaluation is claimed here.

\section{Conclusion}
We proposed a feedback extension to memory-bounded batching: track per-class footprints, preserve a conservative resource bound, and revise future batch capacities when measurements indicate drift. The CPU simulation shows near parity under stable class footprints, where adaptation has little to learn. Whether a safe controller can improve throughput under realistic within-class drift without increasing memory failures remains an open question.

\bibliography{references}
\end{document}
